\documentclass{article}
\usepackage{mathtools} 
\usepackage{fontspec}
\usepackage[UTF8]{ctex}
\usepackage{amsthm}
\usepackage{mdframed}
\usepackage{xcolor}
\usepackage{amssymb}
\usepackage{amsmath}

% 定义新的带灰色背景的说明环境 zremark
\newmdtheoremenv[
  backgroundcolor=gray!10,
  % 边框与背景一致，边框线会消失
  linecolor=gray!10
]{zremark}{命题}

\begin{document}

\begin{zremark}
  \begin{align*}
    log_a x << x^k << c^x << x^x (a > 0, k > 0, c > 1) (x \to +\infty)
  \end{align*}
\end{zremark}


\begin{zremark}
  \begin{align*}
    (\ln\, x)^n << \frac{1}{x^p} \, (x \to 0)
  \end{align*}
  其中$n$是正整数。
\end{zremark}

利用洛必达法则，
\begin{align*}
  \lim\limits_{x \to 0} \frac{(\ln\, x)^n}{\frac{1}{x^p}}
   & = \lim\limits_{x \to 0} \frac{(\ln\, x)^n}{x^{-p}}                                                                       \\
   & = \lim\limits_{x \to 0} \frac{n (\ln x)^{n - 1} \frac{1}{x}}{(-p) x^{-p - 1}}                                            \\
   & = - \frac{n}{p} \cdot \lim\limits_{x \to 0} \frac{(\ln x)^{n - 1}}{x^{-p}}                                               \\
   & = - \frac{n}{p} \cdot \lim\limits_{x \to 0} \frac{(n - 1)(\ln x)^{n - 2}\frac{1}{x}}{(-p) x^{-p - 1}}                    \\
   & = (- \frac{n}{p})(- \frac{n - 1}{p}) \cdot \lim\limits_{x \to 0} \frac{(\ln x)^{n - 2}}{x^{-p}}                          \\
   & \vdots                                                                                                                   \\
   & = (- \frac{n}{p})(- \frac{n - 1}{p}) \cdots (- \frac{2}{p}) \cdot \lim\limits_{x \to 0} \frac{\ln x}{x^{-p}}             \\
   & = (- \frac{n}{p})(- \frac{n - 1}{p}) \cdots (- \frac{2}{p}) (- \frac{1}{p}) \cdot \lim\limits_{x \to 0} \frac{1}{x^{-p}} \\
   & = (- \frac{n}{p})(- \frac{n - 1}{p}) \cdots (- \frac{2}{p}) (- \frac{1}{p}) \cdot \lim\limits_{x \to 0} x^p              \\
   & = 0
\end{align*}

\end{document}